% H9-T19: actual tame cyclic-cover monodromy and exact height norm.
% Run JAC1964-20261121; /root/h9_height; GPT-6 Astra Pro.

\subsection*{H9-T19. Geometric generic tame cyclic rational covers add no height directions}

\begin{theorem}[H9-T19]
Let $A/k$ be an abelian variety of even dimension $n=2m\ge4$,
where $k$ is algebraically closed of arbitrary characteristic.
Let $f:\widetilde A\to\mathbf P^1_t$ be a genuine Lefschetz
pencil for a projective embedding, with critical set $\Sigma$.
Choose the coordinate $t$ so that $\infty\notin\Sigma$.
Let $N\ge2$ be prime to the characteristic, let $s$ be
transcendental over $k$, and put
\[
 F=\overline{k(s)},\qquad K=F(t),\qquad L=F(y),
                   \qquad t=y^N+s.
\]
Write $Y/K$ for the original generic pencil fibre after extension
of constants to $F$. All homologically trivial rational
codimension-$m$ cycles on $Y$ and $Y_L$ are admissible. For the
same signed height convention on both sides,
\[
 h_L(z,w)=\frac1N h_K(\operatorname{cor}z,
                            \operatorname{cor}w)
 \quad\text{for all }z,w\in
                 \operatorname{CH}^m(Y_L)_{\mathrm{hom},\mathbf Q}.
                                                               \tag{H9.93}
\]
In particular the entire new norm-zero subgroup is in the height
radical, and
\[
 \operatorname{rad}h_L
       =\operatorname{cor}^{-1}(\operatorname{rad}h_K).
                                                               \tag{H9.94}
\]
Corestriction identifies the two height quotients, with the form
on the new quotient equal to $1/N$ times the old form.

For $m=2$, accepting 032, the height on the entire homologically
trivial group over $L$ is positive semidefinite by H9-T8.
For arbitrary $m$, (H9.93) transfers whichever positivity or
signature statement has been proved for the original axis form;
it does not assert that the sign of an arbitrary special axis is
known in positive characteristic.
\end{theorem}

\begin{proof}
We prove first the precise cohomological assertion underlying
(H9.93). Fix an invertible prime $\ell$ and allow
$\overline{\mathbf Q}_\ell$ coefficients when decomposing deck
characters. Write $q=|\Sigma|$ and let $V$ be the variable part of
$R^{n-1}f_*\mathbf Q_\ell$ on the smooth-fibre open set. The
fixed part is the constant image of $H^{n-1}(A)$; hard Lefschetz
on $A$ makes its cup-product restriction nondegenerate and gives
the orthogonal splitting into the fixed and variable parts.

By JAC-P07, $V$ has positive rank $r$ and every local vanishing
cycle is nonzero. SGA~7~II, XVIII~6.7.2, makes $V$ absolutely
irreducible, including after finite-index restriction. Its
inertia at each $a\in\Sigma$ is one rank-one transvection,
and it is unramified at $\infty$. It is tame everywhere. The
last assertion also holds in characteristic two for this
odd-dimensional Lefschetz fibre; the tame Picard--Lefschetz
formula is the same one used in H9-T6. In what follows the
cover degree $N$, rather than the characteristic, is restricted.

\paragraph{The actual family of covers and the absence of punctual terms.}
Let $S=\mathbf A^1_s-\Sigma$. Over $S$ form the finite flat map
\[
 \pi:\mathbf P^1_y\times S\longrightarrow\mathbf P^1_t\times S,
 \qquad [y_0:y_1]\longmapsto
                 [y_0^N+s y_1^N:y_1^N].
                                                               \tag{H9.95}
\]
Its group is $\mu_N$, and its branch values are $s$ and $\infty$.
Both are smooth pencil values. Set
\[
 \mathscr X=(\widetilde A\times S)
                   \times_{\mathbf P^1_t\times S}
                          (\mathbf P^1_y\times S),
 \qquad p:\mathscr X\longrightarrow S.
\]
The fibres $\mathcal X_s$ are smooth projective $n$-folds.
Over the smooth locus of $f$, this follows from smooth base
change to the smooth surface $\mathbf P^1_y\times S$; over
$\Sigma$, (H9.95) is \'etale because $s\notin\Sigma$.
The family is flat over $S$, for example from its incidence
equation in $A\times\mathbf P^1_y\times S$, and its geometric
fibres are smooth. Hence $p$ is smooth and proper.

Condition~A says exactly that the maps
$R^i f_*\mathbf Q_\ell\to j_*j^*R^i f_*\mathbf Q_\ell$ are
isomorphisms, where $j$ includes the smooth-fibre locus:
SGA~7~II, XVIII~5.3.2--5.3.5, printed pp.~287--288.
Thus there are no extra singular-fibre stalks in these ordinary
cohomology sheaves. In particular $R^n f_*\mathbf Q_\ell$ and
$R^{n-2}f_*\mathbf Q_\ell$ are constant, by weak Lefschetz and
its dual. In degree $n-1$ the sheaf is the middle extension of
the fixed part plus $V$. Proper base change gives the same
description after (H9.95): the map is unramified over the old
singular fibres and only branches where the old sheaves are
lisse. This is the place where condition~A excludes possible
punctual contributions; merely knowing the generic local
system would not suffice.

For a fixed $s\in S$, the Leray spectral sequence for
$f_s:\mathcal X_s\to\mathbf P^1_y$ has only columns $0,1,2$.
The terms of total degree $n$ are
\[
 H^0(\mathbf P^1_y,R^n f_{s*}\mathbf Q_\ell),\quad
 H^1(\mathbf P^1_y,R^{n-1}f_{s*}\mathbf Q_\ell),\quad
 H^2(\mathbf P^1_y,R^{n-2}f_{s*}\mathbf Q_\ell).
\]
The deck group acts trivially on the two outer terms, since
the indicated sheaves are constant and automorphisms of
$\mathbf P^1$ act trivially on its $H^0$ and $H^2$.
The middle term has no incoming or outgoing differential from
$E_2$ on, since either would lie outside columns $0,1,2$.
The fixed summand contributes no middle term, because
$H^1(\mathbf P^1,\mathbf Q_\ell)=0$. It follows, by taking the
exact character projectors on the Leray filtration, that every
nontrivial deck-character part of $H^n(\mathcal X_s)$ is
canonically the corresponding part of the parabolic cohomology
of $\pi_s^*V$. No splitting of the full Leray filtration is
being assumed.

\paragraph{The nontrivial character parts are actual middle convolutions.}
For a nontrivial character of $\mu_N$, choose the associated
nontrivial Kummer character $\chi$ so that its summand of
$\pi_{s*}\overline{\mathbf Q}_\ell$ is
$\mathcal L_\chi(s-t)$ on the complement of $s,\infty$.
Changing the character to its inverse if necessary only
renames the summands. Finite pushforward and projection formula
now identify the nontrivial part just computed with
\[
 W_{\chi,s}
  =H^1\bigl(\mathbf P^1_t,
        j_{s*}(V\otimes\mathcal L_\chi(s-t))\bigr),
                                                               \tag{H9.96}
\]
where $j_s$ omits $\Sigma,s,\infty$. At an old critical point
the cover is unramified, and at a branch point the old $V$ is
lisse. These two stalk calculations verify that the extension
in (H9.96) is precisely $j_{s*}$, without a missing or extra
skyscraper term.

The preceding construction is functorial in $s$. Relative
Leray, proper base change, and the exact character projectors
therefore identify the corresponding lisse summand of
$R^n p_*\overline{\mathbf Q}_\ell$ with the sheaf whose fibres
are (H9.96). Katz's explicit recipe for middle convolution
identifies this sheaf with $MC_\chi(V)$: \emph{Rigid Local
Systems}, \S5.1.6--5.1.7, especially (3)--(4), chapter-local
pp.~3--4. Thus this is a calculation for the actual cover
family (H9.95), not an application of a full-linear-system
hypersurface-monodromy theorem.

\paragraph{Irreducibility and infinite local monodromy.}
Katz's Theorem~2.9.7 makes middle convolution by a nontrivial
Kummer sheaf an automorphism of the relevant irreducible
perverse objects. Here $V$ is irreducible, tame and has rank at
least two; it has at least two finite singularities and is an
object of Katz's category $\mathcal T_\ell$ by \S5.1.1--5.1.3.
Consequently $W_\chi=MC_\chi(V)$ is a nonzero irreducible
middle-extension local system. Katz~\S5.1.4 and Theorem~4.3.11
give this statement and the following local rules in every
characteristic under consideration, including characteristic
zero; the characteristic-$p$ Fourier-transform calculation is
not being used outside its proved scope.

Here is the exact calculation at infinity. By Katz,
Corollary~3.3.6, there is a tame inertia representation $M$ such
that
\[
 V_\infty=M/M^{I_\infty},\qquad
 (W_\chi)_\infty=(M\otimes\chi)/(M\otimes\chi)^{I_\infty},
 \qquad
 \dim M=\sum_{a\in\Sigma}\dim(V_a/V_a^{I_a})=q.
                                                               \tag{H9.97}
\]
The first quotient is the trivial representation of rank $r$.
It follows directly from Jordan normal form that $M$ is
unipotent, has exactly $r$ blocks of size two, and has $q-2r$
blocks of size one. In particular $q\ge2r$. Tensoring by the
nontrivial $\chi$ removes all invariants, so
\[
 \operatorname{rank}W_\chi=q,\qquad
 (W_\chi)_\infty
       =J(\chi,2)^{\oplus r}\oplus
                            J(\chi,1)^{\oplus(q-2r)}.
                                                               \tag{H9.98}
\]
The displayed character is read using Katz's inertia convention;
replacing an inertia generator by its inverse replaces $\chi$
by $\chi^{-1}$ and has no effect on the argument. A positive
power of inertia at infinity is a nontrivial unipotent, because
$r>0$ and the coefficient field has characteristic zero. Hence
the monodromy image of $W_\chi$ is infinite. For completeness,
the same corollary shows that its local monodromy at each
$a\in\Sigma$ is a pseudoreflection with nontrivial eigenvalue
$\chi$: $V_a/V_a^{I_a}$ is trivial of dimension one. No
classification of reflection groups is required below.

An elementary representation argument now suffices. If an
irreducible representation with infinite image had a nonzero
vector fixed by a finite-index subgroup, replace that subgroup
by its finite-index normal core. Its invariant subspace is
nonzero and stable under the full group, so irreducibility
makes the entire representation invariant. The image would
then factor through a finite quotient, a contradiction. Thus
$W_\chi$ has no nonzero finite-orbit vector, for every
$\chi\ne1$. This conclusion does not assert that its
restriction to every finite-index subgroup is irreducible, nor
does it assert that its full monodromy group is orthogonal.

\paragraph{Geometric generic algebraic classes are deck invariant.}
Every rational algebraic cycle on $\mathcal X_F$ is defined over
a finite extension of $k(s)$, because $F=\overline{k(s)}$.
Its class in $H^n(\mathcal X_F,\mathbf Q_\ell(m))$ is therefore
fixed by a finite-index subgroup of the geometric fundamental
group, after shrinking the parameter curve if needed. Purely
inseparable extensions do not change this assertion. The Tate
twist is constant on geometric monodromy since $k$ is
algebraically closed. Every nontrivial deck projection of that
class is also a finite-orbit vector, as the deck action is
defined over $k$ and commutes with monodromy. By the preceding
paragraph all these projections vanish. We have proved
\[
 \operatorname{cl}\operatorname{CH}^m(\mathcal X_F)_{\mathbf Q}
      \subseteq H^n(\mathcal X_F,\mathbf Q_\ell(m))^{\mu_N}.
                                                               \tag{H9.99}
\]
This is a statement about algebraic cohomology classes. It
does not claim that the new Chow group or its norm kernel is
zero.

\paragraph{The height identity uses actual Chow push--pull.}
Let $Q:\mathcal X_F\to\widetilde A_F$ be the finite flat map
of degree $N$. All homologically trivial generic cycles have
admissible closures on both smooth models, by JAC-P07 and
H9-T12(iii), since the cover is unramified over $\Sigma$.
For any such closure $Z$, form the actual rational Chow cycle
\[
       Z_{\mathrm{new}}=Z-\frac1N Q^*Q_*Z.
\]
The finite flat cyclic-quotient identity
$Q^*Q_*Z=\sum_{g\in\mu_N}g^*Z$ makes its cohomology class
exactly the nontrivial deck part of $[Z]$. This identity can
also be checked as a correspondence: the top-dimensional
cycle of the fibre product is the sum of the deck graphs;
ramification introduces no additional top-dimensional
component. Equation (H9.99) gives $[Z_{\mathrm{new}}]=0$.

The generic restriction of $Q_*Z$ is $\operatorname{cor}z$.
Using this particular closure avoids any uniqueness assertion
for cohomological lifts. A different closure differs by a
vertical cycle and gives the same height by admissibility.
Since a globally homologically trivial closure has zero
intersection with every admissible closure, we obtain
\[
 h_L\left(z-\frac1N\operatorname{res}\operatorname{cor}z,w\right)=0.
\]
Kahn's adjunction
$h_L(\operatorname{res}a,w)=h_K(a,\operatorname{cor}w)$ now
proves (H9.93). Equivalently one may apply ordinary projection
formula to the displayed global closures. The same sign
$(-1)^m$, if included in the height convention, occurs on both
sides and does not alter the factor $1/N$.

Corestriction is surjective over $\mathbf Q$, since
$\operatorname{cor}\operatorname{res}=N$. Testing (H9.93)
against all $w$ therefore proves (H9.94) and the quotient
assertion. For $m=2$, H9-T8 over the algebraically closed field
$F$ supplies semipositivity of $h_K$ and its exact old radical;
(H9.93) gives the stated full-group conclusion over $L$.
\end{proof}

\begin{corollary}[H9-T19a: the full parameter space of rational covers]
Keep the pencil and $N$ of H9-T19. Let $T$ be the open subset of
$\operatorname{Hom}_N(\mathbf P^1_y,\mathbf P^1_t)$ consisting
of maps unramified above $\Sigma$. Let
$\Omega=\overline{k(T)}$, and let
$\pi_\eta:\mathbf P^1_\Omega\to\mathbf P^1_\Omega$ be its
geometric generic map. Put $K=\Omega(t)$ and
$L=\Omega(y)$ using this map. Then (H9.93)--(H9.94), including
the full-group semipositivity assertion for an abelian fourfold
under 032, hold for this geometric generic degree-$N$ rational
cover.
\end{corollary}

\begin{proof}
The usual parameter space of pairs of degree-$N$ binary forms,
modulo common scalar and with nonzero resultant, is an open
subset of $\mathbf P^{2N+1}$. Thus it is smooth and irreducible.
Being unramified over the finitely many points of $\Sigma$ is
an open condition. It is nonempty because it contains the
cyclic curve of maps (H9.95). Hence $T$ is smooth and
geometrically connected. Its universal map
\[
 \Pi:\mathbf P^1_y\times T\longrightarrow\mathbf P^1_t\times T
\]
is finite flat of degree $N$. It is proper with finite fibres,
hence finite. Its source and target are smooth and of the
same dimension; the source is Cohen--Macaulay and the target
is regular, so the zero-dimensional-fibre flatness criterion
gives flatness. The degree on every fibre is $N$.

Set
\[
 \mathscr X_T=(\widetilde A\times T)
                  \times_{\mathbf P^1_t\times T}
                          (\mathbf P^1_y\times T),
 \qquad Q_T:\mathscr X_T\to\widetilde A\times T.
\]
The morphism $p_T:\mathscr X_T\to T$ is smooth and proper.
Near an old critical fibre, $\Pi$ is \'etale, so the assertion
follows from the smoothness of $\widetilde A\times T\to T$.
Elsewhere $f$ is smooth, and base change followed by the
smooth map $\mathbf P^1_y\times T\to T$ gives the assertion.
In particular $R^n p_{T*}\mathbf Q_\ell$ is lisse. The finite
flat pullback and proper trace of $Q_T$ are morphisms of these
relative cohomology sheaves, compatible with arbitrary base
change, and satisfy
\[
 Q_{T*}Q_T^*=N\quad\text{on }H^n(\widetilde A)\otimes\mathbf Q_{\ell,T}.
\]
Consequently
\[
 \mathscr E=\ker Q_{T*}
    =\operatorname{im}\left(1-\frac1NQ_T^*Q_{T*}\right)
                                                               \tag{H9.100}
\]
is a lisse direct summand. This construction is on the actual
cover space, not on a hypothetical full hypersurface system.

Restrict (H9.100) to the cyclic curve $S\to T$ of (H9.95).
There the cyclic push--pull identity identifies $\mathscr E$
with the direct sum of all the nontrivial deck-character
sheaves $W_\chi$. By the proof of H9-T19, this restriction
has no nonzero finite-orbit vector. If the monodromy of
$\mathscr E$ on $T$ had a finite-orbit vector, its restriction
to the image of $\pi_1(S)$ would still have a finite orbit.
This is impossible. No surjectivity of
$\pi_1(S)\to\pi_1(T)$ is needed.

An algebraic class on the geometric generic fibre is fixed
by a finite-index subgroup because its representing cycle is
defined over a finite extension of $k(T)$. Its projection to
$\mathscr E$ must therefore vanish. For an actual Chow
closure $Z$ on that fibre, the cycle
$Z-N^{-1}Q_\eta^*Q_{\eta*}Z$ is algebraic and its class lies
in $\mathscr E$, so its class is zero. The same global
intersection and restriction--corestriction proof as in
H9-T19 now gives (H9.93), with no Galois or deck-group
assumption on the generic rational map. Surjectivity of the
norm over $\mathbf Q$ gives (H9.94).
\end{proof}

\paragraph{Scope and comparison with the negative cover example.}
The constants and genericity in this theorem are essential:
the branch point $s$ is the geometric generic point relative
to the originally fixed pencil over $k$. No assertion is made
that every specialized cyclic cover has only invariant
algebraic middle classes. Special fibres can acquire new
classes. In addition the covering curve here is rational.
For a threefold fibre, its Picard variety is isogenous to the
constant $\operatorname{Pic}^0(A)_K$; every point over
$F(y)$ of that constant abelian variety is constant, because
a rational map from $\mathbf P^1_F$ to an abelian variety
extends and is constant. Thus the Lefschetz Picard-height
directions have zero height here. This is consistent with the
strictly negative directions of H9-T17 over an elliptic
covering curve. The proof of (H9.93) itself does not need this
Picard observation or an Albanese decomposition.

Katz's irreducibility and local rules are cited background.
Their application to the actual moving cyclic-cover family,
the no-finite-orbit argument, and the exact height consequence
are the deductions proved here. The primary source was read
directly in the author's corrected manuscript of
\emph{Rigid Local Systems}: Theorem~2.9.7, chapter~2 p.~30
(PDF p.~67); Corollary~3.3.6, chapter~3 pp.~14--16
(PDF pp.~107--109); Theorem~4.3.11, chapter~4 p.~9
(PDF p.~122); and \S5.1.1--5.1.7, chapter~5 pp.~1--4
(PDF pp.~123--126). The exact local formula was also visually
checked on PDF p.~107. No formal verification or claim of
historical originality is made.
